% =====================================================================
% §5 BTM solar alignment
% =====================================================================

\section{State-level BTM solar alignment}
\label{sec:btm}

Why these zones, and not others? The per-zone drift pattern from
§\ref{sec:drift} aligns with what one would expect if 2022$\to$2026
behind-the-meter (BTM) solar buildout were the dominant drift mechanism.
We collapse the 8 zones to 5 states (CT, MA, ME, NH, RI, VT;
Massachusetts's three zones average to one MA state with load-weighted
mean), pull state-level distributed PV capacity for 2026-Q1 from EIA
Form 861M \citep{eia2025distpv}, normalize by 2024 Census population to
get distributed-PV-MW-per-million-people, and plot against per-state
MAPE drift (W2 minus 2022 self-eval).

\paragraph{Result.}
Figure~\ref{fig:btm} shows the scatter. With $n = 6$ states (CT, MA,
ME, NH, RI, VT — Massachusetts collapsed from its three ISO-NE zones)
and a two-sided permutation test (10\,000 permutations of state
labels), we report Spearman $\rho = \PHRHO{}$ with empirical
$p = \PHPVAL{}$.

\begin{figure}[h]
\centering
\includegraphics[width=0.7\linewidth]{btm_solar_correlation.png}
\caption{State-level distributed PV capacity (MW per million population,
EIA Form 861M, 2026-Q1 cross-section) plotted against per-state MAPE
drift (W2 deployment minus 2022 self-eval, percentage points). Marker
size scales with absolute drift. Spearman $\rho$ and the two-sided
permutation $p$-value are annotated. The $n = 5$ caveat is explicit on
the figure: power is too low to claim a hypothesis test result, but the
direction matches the policy hypothesis.}
\label{fig:btm}
\end{figure}

\paragraph{Causal language.} The observed alignment is consistent with
the BTM-solar mechanism but does not establish it. Other co-varying
drift mechanisms in 2022--2026 New England include EV adoption, energy
efficiency, and post-COVID load patterns; we do not separate them in
this work. Cleaner identification (e.g., a synthetic-counterfactual
demand series that removes BTM solar but holds other dynamics fixed) is
beyond scope. Two facts make the BTM mechanism the leading candidate:
(i) the per-zone signature is geographically concentrated in MA + RI,
the two states with the steepest 2022$\to$2026 distributed-PV growth in
New England, and (ii) the trained baseline depends on a weather signal
that BTM solar systematically distorts (daytime grid load is suppressed
by BTM generation that the model's HRRR inputs do not see), while the
weather-agnostic Chronos predictor is comparatively robust
(§\ref{sec:drift}). Both are circumstantial; we report alignment, not
causation.
